\documentclass[10pt,a4paper]{amsart}
\usepackage{amsmath,amssymb,amsthm,graphicx}
\usepackage{verbatim}
\newenvironment{reference}{\comment}{\endcomment}
%\newenvironment{reference}{}{}
\newenvironment{slogan}{\comment}{\endcomment}
\newenvironment{history}{\comment}{\endcomment}

% For commutative diagrams we use Xy-pic
\usepackage[all]{xy}

% We use 2cell for 2-commutative diagrams.
\xyoption{2cell}
% Theorem environments.
%
\theoremstyle{plain}
\newtheorem{theorem}[subsection]{Theorem}
\newtheorem{proposition}[subsection]{Proposition}
\newtheorem{lemma}[subsection]{Lemma}

\theoremstyle{definition}
\newtheorem{definition}[subsection]{Definition}
\newtheorem{example}[subsection]{Example}
\newtheorem{exercise}[subsection]{Exercise}
\newtheorem{situation}[subsection]{Situation}

\theoremstyle{remark}
\newtheorem{remark}[subsection]{Remark}
\newtheorem{remarks}[subsection]{Remarks}

\numberwithin{equation}{subsection}

% Macros
%
\def\lim{\mathop{\mathrm{lim}}\nolimits}
\def\colim{\mathop{\mathrm{colim}}\nolimits}
\def\Spec{\mathop{\mathrm{Spec}}}
\def\Hom{\mathop{\mathrm{Hom}}\nolimits}
\def\Ext{\mathop{\mathrm{Ext}}\nolimits}
\def\SheafHom{\mathop{\mathcal{H}\!\mathit{om}}\nolimits}
\def\SheafExt{\mathop{\mathcal{E}\!\mathit{xt}}\nolimits}
\def\Sch{\mathit{Sch}}
\def\Mor{\mathop{\mathrm{Mor}}\nolimits}
\def\Ob{\mathop{\mathrm{Ob}}\nolimits}
\def\Sh{\mathop{\mathit{Sh}}\nolimits}
\def\NL{\mathop{N\!L}\nolimits}
\def\CH{\mathop{\mathrm{CH}}\nolimits}
\def\proetale{{pro\text{-}\acute{e}tale}}
\def\etale{{\acute{e}tale}}
\def\QCoh{\mathit{QCoh}}
\def\Ker{\mathop{\mathrm{Ker}}}
\def\Im{\mathop{\mathrm{Im}}}
\def\Coker{\mathop{\mathrm{Coker}}}
\def\Coim{\mathop{\mathrm{Coim}}}

% Boxtimes
%
\DeclareMathSymbol{\boxtimes}{\mathbin}{AMSa}{"02}

%
% Macros for moduli stacks/spaces
%
\def\QCohstack{\mathcal{QC}\!\mathit{oh}}
\def\Cohstack{\mathcal{C}\!\mathit{oh}}
\def\Spacesstack{\mathcal{S}\!\mathit{paces}}
\def\Quotfunctor{\mathrm{Quot}}
\def\Hilbfunctor{\mathrm{Hilb}}
\def\Curvesstack{\mathcal{C}\!\mathit{urves}}
\def\Polarizedstack{\mathcal{P}\!\mathit{olarized}}
\def\Complexesstack{\mathcal{C}\!\mathit{omplexes}}
% \Pic is the operator that assigns to X its picard group, usage \Pic(X)
% \Picardstack_{X/B} denotes the Picard stack of X over B
% \Picardfunctor_{X/B} denotes the Picard functor of X over B
\def\Pic{\mathop{\mathrm{Pic}}\nolimits}
\def\Picardstack{\mathcal{P}\!\mathit{ic}}
\def\Picardfunctor{\mathrm{Pic}}
\def\Deformationcategory{\mathcal{D}\!\mathit{ef}}
\usepackage{cleveref}
\crefname{originalequation}{original equation}{original equations}
\Crefname{originalequation}{Original equation}{Original equations}
\makeatletter
\newlabel{equation-equivalence}{{0BFU}{source}{}{sourcepointer-1}{}}
\newlabel{equation-equivalence@cref}{{[originalequation][0][]0BFU}{[1][0][]source}{}{}{}}
\makeatother
\crefname{originalreference}{original reference}{original references}
\Crefname{originalreference}{Original reference}{Original references}
\makeatletter
\newlabel{lemma-equivalence-fibre}{{0BFX}{source}{}{sourcepointer-2}{}}
\newlabel{lemma-equivalence-fibre@cref}{{[originalreference][0][]0BFX}{[1][0][]source}{}{}{}}
\makeatother
\title{Transferring sequences of blowups in points through the equivalence of formal fibres}
\author{The Stacks Project Authors}
\date{stacks-project:resolve.tex:local-snapshot-20261012; GNU FDL 1.2 or later}
\begin{document}
\maketitle
\noindent\textit{Editorial note.} All delivered mathematical text is the author's own, copied from the pinned Stacks Project chapter file resolve.tex; no statement or proof step was rewritten, reordered or removed.\par\medskip
\noindent\textit{Editorial note.} The counted proof is the author's complete argument: the tracking of the point blown up through the equivalence, the matching of closed points using the isomorphism of the fibres, and the reverse construction.\par\medskip
\noindent\textit{Editorial note.} The label lemma-equivalence-sequence-blowups names the counted statement, and the proof environment immediately following it in the source is delivered verbatim as its proof.\par\medskip
\noindent\textit{Editorial note.} The counted statement is carried by the author's own lemma environment in the chapter source.\par\medskip
\noindent\textit{Editorial note.} Pre-open gate: python operations/agent-authoring/tools/precard\_gate.py --source resolve.tex --label lemma-equivalence-sequence-blowups --start 1494 --end 1533 --local-id 28124 returned RESULT=OK (L28124 free, occupancy FREE, public interval clear). Pre-assembly tag self-check: canonical tag resolved and not present in reference/public-items.jsonl.\par\medskip
\noindent\textit{Editorial note.} The statement and proof are copied from the pinned chapter file; the cited inputs are kept as external pointers with their canonical Stacks tags.\par\medskip
\noindent\textit{Allowed standard prerequisites (editorial).}\begin{itemize}
\item The correspondence between schemes over the base and over the localizations. The equivalence of categories described in the equation labelled equivalence. Transfers morphisms and blowing up data between the base and the localizations at the points.
\item Isomorphism of the fibres of corresponding schemes. Same chapter, the lemma that the fibres of corresponding schemes at the corresponding points are isomorphic. Identifies the fibre over one of the points with the fibre over the corresponding point, which is what allows matching closed points under the correspondence.
\end{itemize}\medskip
\ifdefined\hypertarget\hypertarget{sourcepointer-1}{}\fi
\noindent\textit{Source reference.} Original equation 0BFU: sources/stacks-resolve/pinned-original-source.tex, source lines 1438--1438. Prerequisite pointer invoked at proof line(s) 1496; it names a standard fact of the same chapter that the counted argument uses, and every step of the counted proof built on it is retained verbatim in the delivered body.\par\medskip
\ifdefined\hypertarget\hypertarget{sourcepointer-2}{}\fi
\noindent\textit{Source reference.} Original reference 0BFX: sources/stacks-resolve/pinned-original-source.tex, source lines 1484--1484. Prerequisite pointer invoked at proof line(s) 1523; it names a standard fact of the same chapter that the counted argument uses, and every step of the counted proof built on it is retained verbatim in the delivered body.\par\medskip
\begin{lemma}
\label{lemma-equivalence-sequence-blowups}
Let $S, s_i, S_i$ be as in (\ref{equation-equivalence})
and assume $f : X \to S$ corresponds to $g_i : Y_i \to S_i$ under $F$.
Then there exists a factorization
$$
X = Z_m \to Z_{m - 1} \to \ldots \to Z_1 \to Z_0 = S
$$
of $f$ where $Z_{j + 1} \to Z_j$ is the blowing up of $Z_j$ at a closed
point $z_j$ lying over $\{s_1, \ldots, s_n\}$ if and only if for each
$i$ there exists a factorization
$$
Y_i = Z_{i, m_i} \to Z_{i, m_i - 1} \to \ldots \to Z_{i, 1} \to Z_{i, 0} = S_i
$$
of $g_i$ where $Z_{i, j + 1} \to Z_{i, j}$ is the blowing up of $Z_{i, j}$
at a closed point $z_{i, j}$ lying over $s_i$.
\end{lemma}

\begin{proof}
Let's start with a sequence of blowups
$Z_m \to Z_{m - 1} \to \ldots \to Z_1 \to Z_0 = S$.
The first morphism $Z_1 \to S$ is given
by blowing up one of the $s_i$, say $s_1$. Applying $F$
to $Z_1 \to S$ we find a blowup $Z_{1, 1} \to S_1$ at $s_1$
is the blowing up at $s_1$ and otherwise $Z_{i, 0} = S_i$ for $i > 1$.
In the next step, we either blow up one of the $s_i$, $i \geq 2$
on $Z_1$ or we pick a closed point $z_1$ of the fibre of $Z_1 \to S$
over $s_1$. In the first case it is clear what to do and in
the second case we use that $(Z_1)_{s_1} \cong (Z_{1, 1})_{s_1}$
(Lemma \ref{lemma-equivalence-fibre})
to get a closed point $z_{1, 1} \in Z_{1, 1}$ corresponding to $z_1$.
Then we set $Z_{1, 2} \to Z_{1, 1}$ equal to the blowing up
in $z_{1, 1}$. Continuing in this manner we construct the factorizations
of each $g_i$.

\medskip\noindent
Conversely, given sequences of blowups
$Z_{i, m_i} \to Z_{i, m_i - 1} \to \ldots \to Z_{i, 1} \to Z_{i, 0} = S_i$
we construct the sequence of blowing ups of $S$ in exactly the same manner.
\end{proof}

\end{document}
